\documentclass{article}

% Same NeurIPS 2026 template as the proposal. Compile on Overleaf or with pdflatex.
\PassOptionsToPackage{numbers,sort&compress}{natbib}
\usepackage[preprint]{neurips_2026}

\usepackage[utf8]{inputenc}
\usepackage[T1]{fontenc}
\usepackage{hyperref}
\usepackage{url}
\usepackage{booktabs}
\usepackage{amsfonts}
\usepackage{amsmath}
\usepackage{nicefrac}
\usepackage{microtype}
\usepackage{xcolor}
\usepackage{graphicx}
\usepackage{array}
\usepackage{float}
\usepackage{tabularx}
\setlength{\tabcolsep}{4pt}

\hypersetup{hidelinks,pdftitle={LearnLoop Final Report},pdfauthor={Group 72}}

\makeatletter
\renewcommand{\@noticestring}{COMP7608B Final Project Report -- DRAFT}
\makeatother

% Remove all \todo{} markers before submission.
\newcommand{\todo}[1]{\textcolor{red}{\textbf{[TODO: #1]}}}
\newcommand{\code}[1]{\texttt{\small #1}}

\title{[Group 72] LearnLoop\\A Stateful Agentic Platform for Adaptive Learning\\Final Report}

\author{Fong Chi Ching, Ching Nok Yee, Lai Hok Lam Justin, Zhou Yuxin}

\begin{document}

\maketitle

\begin{tabular}{@{}lll@{}}
\toprule
\textbf{Name} & \textbf{Student ID} & \textbf{Email} \\
\midrule
Fong Chi Ching & 3035299466 & u3529946@connect.hku.hk \\
Ching Nok Yee & 3035283572 & u3528357@connect.hku.hk \\
Lai Hok Lam Justin & 3036657081 & u3665708@connect.hku.hk \\
Zhou Yuxin & 3036791693 & u3679169@connect.hku.hk \\
\bottomrule
\end{tabular}

\vspace{4pt}
\noindent\textbf{Submitted version:} \todo{git tag/commit hash of the submitted snapshot, e.g.\ \code{v1.0-final}; the current working tree has uncommitted changes}.
\textbf{Repository:} \todo{link, if shared}.

%=====================================================================
\section{Summary}

LearnLoop turns a student's uploaded course documents into a source-grounded, adaptive practice loop. Three agents share one durable runtime: the \textbf{Expert} builds a reviewed concept graph, lessons and private question banks; the \textbf{Tutor} chooses the next prerequisite-safe learning action from persistent learner evidence; the \textbf{Assessing Agent} (``Reviewer'' in the UI) grades responses against private keys and rubrics. A deterministic runtime, not the model, decides what is committed to the learner record. We delivered a runnable local application (FastAPI backend, separate durable worker, Next.js frontend, SQLite with sqlite-vec), a prompt/agent Workbench, and an isolated evaluation harness with LLM pseudo-students. Software and failure tests pass for the core loop; live pseudo-student runs are exploratory and do not yet show learning gains or a baseline comparison (Section~\ref{sec:eval}).

%=====================================================================
\section{Users and proposal continuity}

\paragraph{Users and needs.} University students who receive slides and readings but do not know where to start, what depends on what, or which gaps they still have. They need (i) a structured path through \emph{their own} course material, (ii) practice that matches what they currently know, (iii) feedback that is faithful to the source, and (iv) progress that persists across sessions and is not inflated by lucky guesses or hints.

\paragraph{Core requirements and scope.} One course unit of roughly 10--20 concepts from uploaded PDFs or text; every lesson, question and relationship traceable to source passages; demonstration requires two independent rubric-meeting answers to distinct questions at the target difficulty; durable recovery from interrupted work. Scope is a local, single-owner prototype. OCR, multi-user accounts and external web enrichment are out of scope.

\paragraph{Planned vs.\ delivered.} Table~\ref{tab:continuity} compares the proposal with the submitted system.

\begin{table}[H]
\centering
\small
\caption{Proposal continuity. ``Changed'' items keep the objective but alter the means.}
\label{tab:continuity}
\begin{tabularx}{\linewidth}{@{}p{3.6cm}p{1.6cm}X@{}}
\toprule
\textbf{Proposal item} & \textbf{Status} & \textbf{Delivered / reason for change} \\
\midrule
Three agents in one construct--teach--assess loop & Delivered & Separate packages, prompts and durable workflow IDs (\code{course\_preparation}, \code{next\_activity}, \code{assessment}). \\
Source-grounded graph, 3 difficulty levels & Delivered, changed & One scoped primary lesson per objective plus a Basic starter bank; Intermediate/Advanced questions are generated and verified on demand. Reduced setup cost and duplicate teaching across levels. \\
Utility-ranked Tutor planning, versioned plans & Delivered & Runtime ranks eligible candidates; the model picks one of at most 12 and gives a rationale; immutable plan revisions. \\
Rubric assessment, misconceptions, timing, contradictions & Delivered, extended & Four formats (single/multiple answer, fill-in, open) instead of MCQ only. \\
Learner model and transition rules (Sec.~4.2) & Delivered & \code{core/domain.py}; confirmed weakness uses an explicit threshold of 2 failures. \\
Versioned knowledge-base publication & Delivered & Immutable course releases; sessions are pinned to a release. \\
LangGraph orchestration & Changed & Custom Python workflow registry with SQL job leases and checkpoints. \todo{one sentence on why, e.g.\ checkpoint semantics already needed SQL transactions}. \\
PostgreSQL + Qdrant + Docker Compose & Changed & SQLite + sqlite-vec for zero-install teammate setup; PostgreSQL/Qdrant adapter retained as legacy. No Docker. \\
27 prompt configurations, 4 pseudo-student profiles, fixed-sequence and memory-free baselines & Harness delivered; study \textbf{not} run & Live runs so far use one configuration and no live baselines (Section~\ref{sec:eval}). \\
Human review of every flagged item & Not done & 18 review records pending. \todo{either complete them or record the planned replacement by automated verification as a protocol change}. \\
Token/cost/latency traces & Partial & Full per-call ledgers in evaluation and Workbench runs; ordinary app runs do not record usage or enforce aggregate budgets. \\
Shared course with many learners; missions and rewards & Not delivered & Dropped to focus on the learning loop. \\
\bottomrule
\end{tabularx}
\end{table}

\todo{State whether any change above counts as a major objective change under the course change procedure, and cite the approval/notice if one was filed.}

%=====================================================================
\section{Design and implementation}

\subsection{Workflow}

\begin{enumerate}
\item \textbf{Prepare (Expert).} Upload $\rightarrow$ page-aware passages $\rightarrow$ bge-m3 embeddings in sqlite-vec $\rightarrow$ section-by-section goal discovery over \emph{every} passage $\rightarrow$ semantic merge $\rightarrow$ Course Organizer and Relationship Analyst (parallel, independent contexts) $\rightarrow$ Relationship Reviewer audits prerequisites against sources $\rightarrow$ curriculum pyramid $\rightarrow$ lesson, worked example and two Basic questions per objective $\rightarrow$ grounding review and independent key solving $\rightarrow$ atomic release.
\item \textbf{Plan (Tutor).} Load the pinned release and learner evidence; filter by prerequisite readiness; rank candidates; the model chooses one action and difficulty; reuse an unused verified question, otherwise request one from the Expert.
\item \textbf{Assess (Reviewer).} Grade against the private key/rubric; the model adds criterion coverage, quoted misconceptions, contradictions and uncertainty; \code{domain.assess\_evidence()} decides the permitted transition; assessment, evidence, learner version and event commit in one transaction.
\item \textbf{Observe.} The committed state is the next Tutor observation. The learner explicitly clicks \emph{Next}; assessment never silently advances the session.
\end{enumerate}

\todo{Insert the architecture figure (updated from proposal Appendix A; remove LangGraph/Qdrant/Docker labels).}

\subsection{Autonomous decisions and their guards}

\begin{table}[H]
\centering
\small
\caption{Decisions made by agents and how the runtime bounds them.}
\begin{tabularx}{\linewidth}{@{}p{2.8cm}p{1.6cm}X@{}}
\toprule
\textbf{Decision} & \textbf{Agent} & \textbf{Guard / stopping rule} \\
\midrule
Which learning goals a section teaches; how goals merge & Expert & Source IDs must resolve; separate grounding review; at most 2 corrective regenerations, then \code{needs\_review}. \\
Prerequisite vs.\ related edges & Expert (3 sub-roles) & Reviewer may delete/downgrade; deterministic acyclicity, coverage and known-ID checks. \\
Is a drafted question correct? & Expert verifier & Fresh call solves it without the key; disagreement triggers repair or second verifier; a second opinion cannot approve a rejected draft. \\
Next action, concept, difficulty & Tutor & Chooses only from runtime-validated candidates; invented IDs rejected; prerequisite and version rechecked at commit. \\
Request a targeted diagnostic question & Tutor & Private to learner; falls back to a saved verified question if verification fails. \\
Coverage, misconception, uncertainty & Reviewer & Quotes must appear in the learner's text; cannot turn a wrong choice into a correct one; partial, hinted or uncertain answers never establish demonstration. \\
Clarify vs.\ continue & Runtime & Missing credentials pause the run; unsupported evidence stops for review; learner controls \emph{Next}. \\
\bottomrule
\end{tabularx}
\end{table}

\subsection{Models, prompts, tools, data and state}

\textbf{Models.} Generation: \code{openai/gpt-6-luna} via OpenRouter for all three agents (configurable through \code{OPENROUTER\_MODEL}); optional \code{EXPERT\_VERIFIER\_MODEL} for second reviews. Production calls use the provider's default temperature; evaluation runs fix temperature 0.2 and a seed. Embeddings: BAAI \code{bge-m3} (1024-d) served by Ollama. Resource screening additionally used \code{google/gemini-2.5-flash}. \todo{Record the provider-resolved model version string from a usage ledger.}

\textbf{Prompts.} One file per agent: \code{agents/expert/prompts.py}, \code{agents/tutoring/prompts.py}, \code{agents/assessing/prompts.py}. The Workbench lets an expert user save prompt templates per agent (migration \code{0007}); evaluation prompt variants live in \code{evaluation/variants.py}. All outputs are schema-constrained JSON validated by Pydantic (\code{core/schemas.py}).

\textbf{Tools.} Application operations executed by Python and visible in the activity trace: scoped retrieval (\code{repository.retrieve}), grounding review, publish release, load learner state, assign activity, assess response, record assessment, revise plan, request Expert question. This is a bounded workflow with model-assisted decisions, not an open-ended ReAct tool loop.

\textbf{Data and state.} SQLite (WAL, \code{BEGIN IMMEDIATE} writes) holds documents, passages, vectors, immutable releases, sessions, activities, per-difficulty learner evidence, plans, runs, jobs and validated public events. Answer keys and rubrics never leave the server before assessment.

\subsection{Key choices and trade-offs}
\begin{itemize}
\item \textbf{Deterministic commit rules over model judgment} for mastery: safer and testable, but the transition thresholds are hand-set rather than learned.
\item \textbf{Exhaustive section review} instead of top-$k$ sampling: complete coverage, but preparation is slow (Section~\ref{sec:nfr}).
\item \textbf{Starter bank + on-demand generation}: cheaper setup, but first-time Intermediate/Advanced questions add latency during practice.
\item \textbf{SQLite + sequential worker}: trivial setup and strong consistency; no parallel job lanes and single-owner only.
\end{itemize}

\subsection{Task traces}
\todo{Add 2--3 short trace excerpts (run ID, ordered events, tool, decision, committed change) exported from Agent activity: (a) an ordinary practice turn, (b) a changed request, e.g.\ the learner switches to a blocked topic and receives the \code{unmet\_prerequisites} 409 with a suggested prerequisite, (c) a failure case from Section~\ref{sec:rel}. Link each step to the function that produced it.}

%=====================================================================
\section{Reliability and recovery}
\label{sec:rel}

Each case below is covered by an automated test with fake model/embedding providers on temporary SQLite files.

\begin{table}[H]
\centering
\small
\caption{Tested failure situations.}
\begin{tabularx}{\linewidth}{@{}p{2.4cm}XXp{2.6cm}@{}}
\toprule
\textbf{Failure} & \textbf{Detection and bounded response} & \textbf{User feedback} & \textbf{Test} \\
\midrule
Model/provider error during preparation & Worker records the error; up to 6 attempts; retry receives the failure diagnostics; completed sections and packs are reused. & Agent activity shows ``Work interrupted'' with Retry/Resume. & \code{test\_section\_review\_resume\_only\_reprocesses\_failed\_section}, \code{test\_course\_preparation\_retry\_sends\_failure\_feedback\_to\_expert} \\
Unsupported generated content & Grounding review rejects; 2 repairs with widening evidence (6/12/18 passages); then withheld as \code{needs\_review}. & Document review lists source issues without leaking keys. & \code{test\_semantic\_failure\_blocks\_graph\_publication}, \code{test\_existing\_failed\_pack\_is\_repaired\_or\_remains\_withheld} \\
Duplicate or stale answer & \code{operationId} replay returns the original receipt; conflicting replay or stale \code{expectedLearnerVersion} returns 409. & Client reloads current state. & \code{test\_stale\_version\_rejected...}, \code{test\_assessment\_waits\_for\_explicit\_next\_and\_retry\_keeps\_evidence\_once} \\
Crash during assessment commit & Single transaction: evidence, version and events roll back together. & Run retryable; no partial progress. & \code{test\_assessment\_commit\_failure\_rolls\_back...} \\
Worker crash / expired lease & File lock and lease reclaim; no duplicate chunks; no extra attempt consumed while another worker holds the lock. & Job resumes automatically. & \code{test\_expired\_job\_lease\_is\_reclaimed...} \\
Missing credentials & Run set to \code{paused}, no model call. & ``Configure the generation model, then resume''. & \todo{cite test or manual check} \\
Experiment budget exhausted & Call/token reservations stop the episode; checkpoint kept; no phantom metrics. & Marked failed/stopped in report. & \code{test\_budget\_and\_invalid\_action\_keep\_failed\_attempt...} \\
\bottomrule
\end{tabularx}
\end{table}

\paragraph{Live failures observed.} In the exploratory pseudo-student run \code{722a04b0}, the Novice episode failed with a \code{ValueError} after 22 calls; it was retried individually (runs \code{c01c4eef}, \code{1fedf3cb}) without rerunning the other learners. Another run stopped on a provider \code{RemoteProtocolError}; usage for that attempt is marked incomplete. \todo{root cause of the Novice ValueError}.

\paragraph{Risks from repeated external actions.} Retries are paid model calls. Provider retries (408/425/429/5xx, 3 retries with backoff) count against call/token limits, and an attempt that disconnects after the provider may have processed it is recorded with incomplete usage. Learner-facing writes are idempotent by operation ID, so a retry cannot double-count evidence. Ordinary app runs have per-request limits but no aggregate cost cap.

\todo{The current working tree fails 6 of 265 backend tests, including \code{test\_existing\_failed\_pack\_is\_repaired\_or\_remains\_withheld[unrepairable]} (an unrepairable pack now completes instead of stopping for review). Fix before submission, otherwise the ``unsupported content is withheld'' claim does not hold.}

%=====================================================================
\section{Non-functional requirements}
\label{sec:nfr}

\begin{table}[H]
\centering
\small
\caption{Targets, measures and observed results. Latencies come from live pseudo-student run \code{722a04b0} (3 completed episodes) unless stated.}
\begin{tabularx}{\linewidth}{@{}p{1.7cm}p{2.6cm}X@{}}
\toprule
\textbf{Quality} & \textbf{Target} & \textbf{Measures and result} \\
\midrule
Latency & \todo{e.g.\ Tutor/Reviewer p95 $<$ 10\,s per step} & Tutor median 3.9--4.6\,s, p95 up to 8.6\,s; Reviewer median 4.1--4.6\,s, p95 up to 7.4\,s. Course preparation for 5 concepts: 29 calls, 652\,s. Model calls run outside DB write locks; 90\,s request timeout. \\
Cost & \todo{e.g.\ $<$ US\$0.01 per tutoring episode} & Episode of 3 activities: about US\$0.007 (about 32k input / 7k output tokens, including simulator calls). Preparation of 5 concepts: US\$0.041 (87k / 65k tokens). Starter bank avoids generating unused levels. \\
Workload limits & Bounded inputs & 20\,MiB per file, 200 PDF pages, 1000 chunks; 6 attempts per run; evaluation caps of 80 calls / 400k tokens per prepared concept and 64 calls / 250k tokens per episode. \\
Privacy & No secrets or keys exposed & Bearer token on all \code{/api/*}; owner derived server-side; answer keys/rubrics private until assessed; credentials only in \code{.env} (never \code{NEXT\_PUBLIC\_}); course text is sent to OpenRouter (disclosed assumption). \\
Correctness / integrity & No partial or stale writes & Version checks, operation IDs, atomic commits, immutable hashed releases. \\
\bottomrule
\end{tabularx}
\end{table}

\paragraph{Operating assumptions.} One local user; Ollama with \code{bge-m3} reachable; OpenRouter credit; text-based PDFs (no OCR). \todo{Measure ordinary app latency for one practice turn outside the simulator (the figures above include simulator overhead in end-to-end time).}

%=====================================================================
\section{Evaluation}
\label{sec:eval}

\subsection{Software and acceptance tests}
The backend suite has \todo{N} tests \todo{re-run on the submitted version; currently 259 pass and 6 fail}. A live acceptance journey (19 Oct, real model, embeddings and worker on an isolated database) passed all 21 conditions, including hint-assisted evidence, two independent checks before demonstration, duplicate-answer protection and cross-session continuation.

\subsection{Controlled pseudo-student harness}
\textbf{Design.} Four hidden profiles (novice, intermediate, misconception-prone, hint-dependent) defined in \code{evaluation/simulator.py}; the Tutor and Reviewer never see hidden state. Held-out pre/post MCQ exams, one item per concept per level, scored independently of the experimental Reviewer. Conditions: adaptive (full system), fixed-sequence, and memory-free baselines; 27 prompt combinations (3 Expert $\times$ 3 Tutor $\times$ 3 Assessor emphases).
\textbf{Metrics.} Learning gain and gain per activity; post-test mastered concepts; learner-state macro F1 and difficulty-state accuracy against hidden state; false-demonstration rate; misconception precision/recall; prerequisite violations, invalid updates, repeated questions; tokens, cost, median/p95 latency per role.
\textbf{Settings.} \code{openai/gpt-6-luna}, temperature 0.2, recorded seed; manifests hash code, prompts and resources. Episodes with $<75\%$ agreement between simulated answers and the profile's controlled ability draws are excluded as \code{simulation\_invalid}.

\subsection{Results}

\begin{table}[H]
\centering
\small
\caption{Exploratory live run \code{722a04b0} (16 Nov 2026): 5 concepts from the ``Large Language Model'' course release v6, configuration grounded/utility/weighted, adaptive only, 3 activities, 1 repeat. Not research evidence.}
\label{tab:explore}
\begin{tabular}{@{}lccccccc@{}}
\toprule
Profile & Pre & Post & Gain & Mastered & Macro F1 & Diff.\ acc. & Cost (US\$) \\
\midrule
Novice & \multicolumn{7}{c}{failed (\code{ValueError}); retry: 0.27 $\rightarrow$ 0.00, F1 0.00} \\
Intermediate & 0.67 & 0.33 & $-0.33$ & 2/5 & 0.14 & 0.40 & 0.0066 \\
Misconception-prone & 0.27 & 0.20 & $-0.07$ & 3/5 & 0.20 & 0.67 & 0.0072 \\
Hint-dependent & 0.27 & 0.47 & $+0.20$ & 3/5 & 0.13 & 0.33 & 0.0074 \\
\bottomrule
\end{tabular}
\end{table}

Across completed episodes there were 0 prerequisite violations, 0 invalid state updates and 0 repeated questions, and simulator policy agreement was 97--100\%. Misconception precision was 0.0 where measurable. The 3 Nov one-concept live canary also exposed a failure: the learner was estimated \emph{demonstrated} while the post-test fell from 0.67 to 0.33 (false-demonstration rate 1.0). An offline run of all 27 configurations $\times$ 4 profiles $\times$ 2 repeats (216 episodes) plus 24 baseline episodes completed without workflow failures, but uses fixture answers and only proves the harness runs.

\paragraph{Interpretation.} With three activities per episode and one held-out item per concept-level, gains are dominated by noise; the runs check workflow behaviour, not learning effect. Low macro F1 is expected when most concepts remain unassessed after three activities, but it also means LearnLoop's learner estimate cannot yet be claimed accurate.

\todo{MISSING for the rubric: (1) a \textbf{fair baseline comparison} -- run fixed-sequence and memory-free against the same frozen materials, profiles and activity budget (\code{--baselines}); (2) \textbf{repeats} ($\geq$3) and report mean $\pm$ SD or CI; (3) longer episodes (e.g.\ 24 activities) so states can move; (4) \textbf{human check of LLM judging} -- a team member reviews a sample of grounding/assessment judgments against explicit criteria and reports agreement; (5) a short list of representative, challenging and failure cases with reference outcomes.}

\paragraph{Limitations.} The simulator's ability, learning and forgetting parameters are assumptions, not a validated human model. One course unit; one generation model; no real-learner study.

%=====================================================================
\section{Contributions and attribution}

\begin{table}[H]
\centering
\small
\begin{tabularx}{\linewidth}{@{}p{3cm}X@{}}
\toprule
\textbf{Member} & \textbf{Contribution} \\
\midrule
Fong Chi Ching & \todo{} \\
Ching Nok Yee & \todo{} \\
Lai Hok Lam Justin & \todo{} \\
Zhou Yuxin & \todo{} \\
\bottomrule
\end{tabularx}
\end{table}

\paragraph{Reused software and data.} Backend: FastAPI 0.128.4, Uvicorn 0.40.0, SQLAlchemy 2.0.46, Alembic 1.20.0, Pydantic 2.12.5, pydantic-settings 2.12.0, sqlite-vec 0.1.9, httpx 0.28.1, pypdf 6.11.0, filelock 4.2.8, jsonschema 4.26.0. Frontend: Next.js 16.4.2, React 19.3.0, zod, react-markdown, KaTeX. Services: OpenRouter, Ollama, BAAI bge-m3; optional Poppler \code{pdftoppm}. Data: COMP7608B lecture notes (Lecture 2 for the evaluation unit) \todo{and the source of \code{test.pdf}}. Ideas: ReAct, Reflexion, knowledge tracing, RAG (proposal references).

\paragraph{Unfinished or mocked functionality.} Tests and offline evaluation use fake model/embedding providers and fixture answers (labelled \code{workflow\_only}). Not implemented: LangGraph, Docker, multi-user accounts, missions/rewards, OCR, Assessor-initiated question requests, production-wide usage accounting, cancellation of runs, staff approval of review-required content.

\section*{AI-use declaration}
\todo{List AI tools/models used (e.g.\ Cursor agent with model names, ChatGPT) and what each helped with (code, docs, report drafting). Attach project prompts in the ZIP, redacting credentials and personal data; declare any missing prompt records.}

\section*{References}
\small
[1] Yao et al.\ (2023). ReAct. ICLR. \quad
[2] Ma et al.\ (2014). Intelligent tutoring systems and learning outcomes. J.\ Educ.\ Psych.\ 106(4). \quad
[4] Lewis et al.\ (2020). Retrieval-augmented generation. NeurIPS. \quad
[5] Corbett \& Anderson (1994). Knowledge tracing. UMUAI 4. \quad
[6] Shinn et al.\ (2023). Reflexion. NeurIPS. \quad
[10] BAAI bge-m3 model card. \todo{renumber and complete from proposal reference list}

\end{document}
